\section{Methods}
\subsection{Clustering algorithms}
\subsubsection{PAM} %? da togliere?
The PAM algorithm is very similar to \(k\)-means, mostly because both are partitional algorithms, in other words, both break the dataset into groups (clusters), and both work by trying to minimize the error, but PAM works with Medoids, that are an entity of the dataset that represent the group in which it is inserted, and \(k\)-means works with Centroids, that are artificially created entity that represent its cluster.

The PAM algorithm partitions the dataset of \(n\) objects into \(k\) clusters, where both the dataset and the number \(k\) is an input of the algorithm. This algorithm works with a matrix of dissimilarity, whose goal is to minimize the overall dissimilarity between the representants of each cluster and its members. The algorithm uses the following model to solve the problem:
\[F(x) = \text{minimize} \sum_{i=1}^{n} \sum_{j=1}^{n} d(i,j) z_{ij}\]
Subject to:\\
\(\sum_{i=1}^n z_{ij} = 1, j=1,2,..,n\)\\
\(z_{ij} \leq y_{i}, i,j = 1,2,..,n\)\\
\(\sum_{i=1}^n y_{i} = k, k = \text{number of clusters}\)\\
\(y_{i}, z_{ij} \in \{0,1\} , i,j = 1,2,..,n\)\\
Where:
\begin{itemize}
    \item \(d(i,j)\): dissimilarity measurement between the entities \(i\) and \(j\)
    \item \(z_{ij}\): a variable that ensures that only the dissimilarity between entities from the same cluster will be computed in the main function
\end{itemize}
The others expressions are constraints that have the following functions: 
\begin{enumerate}
    \item ensures that every single entity is assigned to one cluster and only one cluster, 
    \item ensures that the entity is assigned to its medoid that represents the cluster, 
    \item ensures that there are exactly k clusters and 
    \item lets the decision variables assume just the values of 0 or 1. 
\end{enumerate}
2 kinds of input:
\begin{itemize}
    \item matrix of every entity and values of its variables
    \item dissimilarity matrix
\end{itemize}
\paragraph{Algorithm}
Input:\\
\(O\): set of objects\\
\(S\): set of selected objects, medoids\\
\(U = O - S\): set of unselected objects\\
\(D_p\): dissimilarity between \(p\) and the closest object in \(S\), \(p \in S\) if and only if \(D_p = 0\)\\
\(E_p\): the dissimilarity between \(p\) and the second closest object in \(S\), \(D_j \leq E_j\)
\begin{itemize}
    \item \textbf{BUILD PHASE}
    \begin{enumerate}
        \item Choose \(k\) entities to become the medoids (set \(S\)), or in case these entities were provided use them as the medoids;
        \item Calculate the dissimilarity matrix if it was not informed;
        \item Assign every entity to its closest medoid;
    \end{enumerate}
    \item \textbf{SWAP PHASE}
    \begin{enumerate}
        \setcounter{enumi}{3}
        \item For each cluster search if any of the entities of the cluster (elemnt of set \(U\)) lower the average dissimilarity coefficient, if it does select the entity that lowers this coefficient the most as the medoid for this cluster;
        \item If at least one medoid (element of \(S\)) has changed go to (3), else end the algorithm.
    \end{enumerate}
\end{itemize}
\subsubsection{CLARA} %? da togliere?
\subsubsection{CLARANS}
A way to reduce the number of distances calculated to create a clustering could be the use of the CLARANS algorithm, name that stands for Clustering Large Applications based on RANdomized search.

The CLARANS algorithm is different from PAM and CLARA: it is not performed a check for every pair (medoid, non-medoid) to find a possible diminishing of the loss function (sum of square distances from the nearest medoid), instead it is thought an underlying graph abstraction of the problem, where each node is a collection of k points, representing the medoids and every arch is the swap of a medoid with a non-medoid.

The CLARANS algorithm searches in this highly connected graph for some neighbours that bring a lowering of the loss function - and it does so until it finds a local minimum. For each step, the number of neighbours checked is not the entire set of nodes at distance 1 from the current one, instead there is a parameter, \textit{numneighbours}, that manages the level of time spent on a node.

Moreover, after finding a certain amount - numlocal - of local minima, the one that leads to the lowest value of the loss function is returned

We have to point out that, although this algorithm does not find exact local minima, but accepts of calling local minimum a node where the neighbours explored do not better the loss function, the results shown by THIS PAPER [
%todo trova qualcosa sulla qualità]
] show that CLARANS does not miss good solutions in terms of the PAM one, which instead searches for all the neighbours of a set of medoids.

Adding to this, it is clear that being the graph higly connected, moving to the real local minimum does not suffer highly from "taking the wrong direction", because there are multiple routes that bring to the same optimal node.

The pseudocode for CLARANS is the following:

CO

DE

CO

DE


\cite{CLARANS}
\subsection{Distances}
Having the goal of clustering metagenomic datasets, it is necessary to find a way to measure some sort of distance between two sets of reads from possibly different species.

A way could be to compare the reads (or contigs) to taxonomic profiles to find the sample composition in terms of species.

Another way is instead the evaluation of two samples starting from the raw data, the reads, comparing the kmers in common and not in common.
From a read of m bases, m-k+1 kmers of length k are exctracted and put in a set. This can be done for each read in a sample, resulting in a set for each sample.

It is then possible to calculate on these sets some ecological distances, both abundance-based or presence-based. The former returns an higher grade of similarity if a kmer is present in the two samples for which the distance is calculated contain that kmer and they have lots of copies of it inside them.
Presence-based distances instead focus just on the presence/absence of a kmer in both sets, without considering the absolute value of their presence.

Presence-based methods are good if functional/taxonomical patterns are searched, for which even just the presence of a specific DNA sequence is enough to infer some conclusion.

For this type of distance calculation, it is preferred to use abundanc-based distances, because kmers are too small to contain functional information, and at the same time abundance can give information about a higher "overlap rate".

A tool that implements the kmer-count distance definition is Simka [CITARE EL PAPER todo]. From the paper of presentation of this method and realative overview of the Simka algorithm, the distances metrics used in this report have been extrapolated. They are presented in the next sections.
\subsubsection{Bray-Curtis}
It is defined as:
\begin{equation}
    BrayCurtis_{Ab}(S_i, S_j) = 1 - 2\frac{\sum_{w \in S_i\cap S_j}min(N_{S_i}(w), N_{S_j}(w))}{\sum_{w \in S_i}N_{S_i}(w) + \sum_{w \in S_j}N_{S_j}(w)}
\end{equation}
where $w$ is a k-mer and $N_{S_i}(w)$ is the abundance of $w$ in the dataset $S_i$.
\subsubsection{Jaccard}

\subsection{Evaluation metrics} %todo add citation
Evaluating how well the clustering has performed can be approached as a classification problem. In a typical binary classification setting, predicted values are compared with the true (actual) values, and each prediction can be categorized as either Positive or Negative.

When evaluating classification results, there are four possible outcomes:
\begin{itemize}
\item TP (True Positive): The predicted value is Positive, and the actual value is also Positive.
\item FP (False Positive): The predicted value is Positive, but the actual value is Negative.
\item FN (False Negative): The predicted value is Negative, but the actual value is Positive.
\item TN (True Negative): The predicted value is Negative, and the actual value is also Negative.
\end{itemize}
Two commonly used metrics are:
\begin{itemize}
    \item Precision, which measures how many of the predicted Positive values are actually Positive: TP/(TP+FP)
    \item Recall, which measures how many of the actual Positive values are correctly predicted: TP/(TP+FN)
    %\item Accuracy, which measures the proportion of correct classification, both TP and TN: (TP+TN)/(TP+TN+FP+FN)
    \item F1-score, the harmonic mean of precision and recall: 2 \(\times\) (precision \(\times\) recall)/(precision + recall)
\end{itemize}

But in this project we are dealing with a multi-class classification task, we need to categorize each sample into 1 of \(N\) different classes. We need to combine the per-class scores into a single number, the classifier's overall score.

Define:
\begin{itemize}
    \item \(N\): total number of classes
    \item \(n\): total number of actual values
    \item \(x_\alpha\): number of actual values for the class \(\alpha\)
    \item \(P_\alpha\): precision for the class \(\alpha\)%, computed as \(P_A = \frac{TP_A}{TP_A + FP_A}\)
    \item \(R_\alpha\): precision for the class \(\alpha\)%, computed as \(R_A = \frac{TP_A}{TP_A + FN_A}\)
    \item \(F1_\alpha\): F1-score for the class \(\alpha\)
\end{itemize} 

\paragraph{Macro-averaged scores}
Each score is computed as a simple arithmetic mean of our per-class scores.
\begin{equation}
    P_\text{macro-averaging} = \frac{P_A + P_B + \dots + P_N}{N}
\end{equation}
\begin{equation}
    R_\text{macro-averaging} = \frac{R_A + R_B + \dots + R_N}{N}
\end{equation}
\begin{equation}
    F1_\text{macro-averaging} = \frac{F1_A + F1_B + \dots + F1_N}{N}
\end{equation}
\paragraph{Weighted-averaged scores}
The scores of each class are weighted by the number of samples from that class, differently from macro-averaged where each class has same weigth.
\begin{equation}
    P_\text{weighted-averaging} = \frac{P_A x_A + P_B x_B + \dots + P_N x_N}{n}
\end{equation}
\begin{equation}
    R_\text{weighted-averaging} = \frac{R_A x_A + R_B x_B + \dots + R_N x_N}{n}
\end{equation}
\begin{equation}
    F1_\text{weighted-averaging} = \frac{F1_A x_A + F1_B x_B + \dots + F1_N x_N}{n}
\end{equation}
\paragraph{Micro-averaged scores}
Gives equal weight to each instance. We first compute micro-averaged precision and micro-averaged recall over all the samples, and then combine the two.
\begin{equation}
    P_\text{micro-averaging} = \frac{TP_A + TP_B + \dots + TP_N}{TP_A + FP_A + TP_B + FP_B + \dots + TP_A + FP_N}
\end{equation}
\begin{equation}
    R_\text{micro-averaging} = \frac{TP_A + TP_B + \dots + TP_N}{TP_A + FN_A + TP_B + FN_B + \dots + TP_A + FN_N}
\end{equation}
We can observe that \(FP_A + FP_B + \dots + FP_N = FN_A + FN_B + \dots + FN_N\), because are all the values falsely classified. So 
\begin{equation}
    R_\text{micro-averaging}= P_\text{micro-averaging} = F1_\text{micro-averaging}
\end{equation}

%for evaluation we can use macro- or micro-average precision and recall (\cite{MultiClass_Evaluation_Metrics}). Macro-averaging gives equal weight to each class, while micro-averaging gives equal weight to each instance.

% Accuracy $ACC$ measures the proportion of correctly classified cases from the total number of objects in the dataset.
% \begin{equation}
%     ACC = \frac{\text{correct prediction}}{\text{all predictions}}
% \end{equation}

% accuracy has its downsides. While it does provide an estimate of the overall model quality, it disregards class balance and the cost of different errors. 

% You can use them in multi-class classification problems as well, there are different approaches to how to do that, each with its pros and cons. 

% \begin{itemize}
%     \item Precision and recall by class
    
%     Precision $P_A$ for a given class $A$ in multi-class classification is the fraction of instances correctly classified as belonging to a specific class out of all instances the model predicted to belong to that class. 

%     \begin{equation}
%         P_A = \frac{TP_A}{TP_A + FP_A}
%     \end{equation}

%     Recall $R_A$ for a given class $A$ in multi-class classification is the fraction of instances in a class that the model correctly classified out of all instances in that class. 

%     \begin{equation}
%         R_A = \frac{TP_A}{TP_A + FN_A}
%     \end{equation}

%     \item Averaging precision and recall
    
%     \begin{itemize}
%         \item Macro-averaging
%         \begin{equation}
%             P_\text{macro-averaging} = \frac{P_A + P_B + \dots + P_N}{N}
%         \end{equation}
%         \begin{equation}
%             R_\text{macro-averaging} = \frac{R_A + R_B + \dots + R_N}{N}
%         \end{equation}
%         \item Micro-averaging
%         \begin{equation}
%             P_\text{micro-averaging} = \frac{TP_A + TP_B + \dots + TP_N}{TP_A + FP_A + TP_B + FP_B + \dots + TP_A + FP_N}
%         \end{equation}
%         \begin{equation}
%             R_\text{micro-averaging} = \frac{TP_A + TP_B + \dots + TP_N}{TP_A + FN_A + TP_B + FN_B + \dots + TP_A + FN_N}
%         \end{equation}
%     \end{itemize}
    

% \end{itemize}

% In the first case, you can calculate the precision and recall for each class individually. This way, you get multiple metrics based on the number of classes you have in the dataset.\\
% In the second case, you can "average" precision and recall across all the classes to get a single number. You can use different methods to average the precision, such as macro- or micro-averaging.

% Moreover it is possible to use F1 as a parameter that joins the information given by both precision and recall. The formula of F1 for any class A is:
% \begin{equation}
%             F1_A = 2 \cdot \frac{P_A \cdot R_A}{P_A + R_A}
% \end{equation}

% In the context of multiclass evaluation, you can either calculate the F1-score for each class, or instead averaging the different values obtained in an only parameter. As before this can be done with macro-averaging or micro-averaging.

% The formulas for the two types of averaging are:
% \begin{equation}
%             F1_\text{macro-averaging} = \frac{F1_A + F1_B + \cdots + F1_N}{N}
% \end{equation}

% \begin{equation}
%             F1_\text{micro-averaging} = 2 \cdot \frac{P_\text{micro-averaging} \cdot R_\text{micro-averaging}}{P_\text{micro-averaging} + R_\text{micro-averaging}}
% \end{equation}

\subsection{Dataset}
The dataset used for the report consists in one project of the GOS (Global Ocean Sampling) marine dataset, which contains reads sequenced with Sanger (prevalently) and pyrosequencing.

The sample we have chosen are the following, listed with their geographical label, that is where the sample had been taken.

\begin{table}[ht]
\centering
\begin{tabular}{|c|l|c|}
\hline
\textbf{Sample ID} & \textbf{Geographical Label} & \textbf{Dimension (KB)} \\
\hline
GOS02 & North American East Coast & --- \\
GOS03 & North American East Coast & --- \\
GOS04 & North American East Coast & --- \\
GOS05 & North American East Coast & --- \\
GOS06 & North American East Coast & --- \\
GOS07 & North American East Coast & --- \\
GOS08 & North American East Coast & --- \\
GOS09 & North American East Coast & --- \\
GOS10 & North American East Coast & --- \\
GOS11 & North American East Coast & --- \\
GOS12 & North American East Coast & --- \\
GOS13 & North American East Coast & --- \\
GOS14 & North American East Coast & --- \\
GOS15 & Caribbean Sea             & --- \\
GOS16 & Caribbean Sea             & --- \\
GOS17 & Caribbean Sea             & --- \\
GOS18 & Caribbean Sea             & --- \\
GOS19 & Caribbean Sea             & --- \\
GOS20 & Panama Canal              & --- \\
GOS21 & Eastern Tropical Pacific  & --- \\
GOS22 & Eastern Tropical Pacific  & --- \\
GOS23 & Eastern Tropical Pacific  & --- \\
GOS24 & Eastern Tropical Pacific  & --- \\
GOS25 & Galapagos Islands         & --- \\
GOS26 & Galapagos Islands         & --- \\
GOS27 & Galapagos Islands         & --- \\
GOS28 & Galapagos Islands         & --- \\
GOS29 & Galapagos Islands         & --- \\
GOS30 & Galapagos Islands         & --- \\
GOS31 & Galapagos Islands         & --- \\
GOS32 & Galapagos Islands         & --- \\
GOS33 & Galapagos Islands         & --- \\
GOS34 & Galapagos Islands         & --- \\
GOS35 & Galapagos Islands         & --- \\
GOS36 & Eastern Tropical Pacific  & --- \\
\hline
\end{tabular}
\caption{Sample metadata: Part 1 of 2}
\label{tab:samples_part1}
\end{table}
\begin{table}[ht]
\centering
\begin{tabular}{|c|l|c|}
\hline
\textbf{Sample ID} & \textbf{Geographical Label} & \textbf{Dimension (KB)} \\
\hline
GOS46 & Tropical South Pacific      & --- \\
GOS47 & Polynesia Archipelagos     & --- \\
GOS48 & Polynesia Archipelagos     & --- \\
GOS49 & Polynesia Archipelagos     & --- \\
GOS50 & Polynesia Archipelagos     & --- \\
GOS51 & North American East Coast  & --- \\
GOS52 & Polynesia Archipelagos     & --- \\
GOS53 & Indian Ocean               & --- \\
GOS54 & Indian Ocean               & --- \\
GOS55 & Indian Ocean               & --- \\
GOS56 & Indian Ocean               & --- \\
GOS57 & Indian Ocean               & --- \\
GOS59 & Indian Ocean               & --- \\
GOS60 & Indian Ocean               & --- \\
GOS61 & Indian Ocean               & --- \\
GOS62 & Indian Ocean               & --- \\
GOS63 & Indian Ocean               & --- \\
GOS64 & Indian Ocean               & --- \\
GOS65 & Indian Ocean               & --- \\
GOS66 & Indian Ocean               & --- \\
GOS67 & Indian Ocean               & --- \\
GOS68 & Indian Ocean               & --- \\
GOS69 & Indian Ocean               & --- \\
GOS70 & Indian Ocean               & --- \\
GOS71 & Indian Ocean               & --- \\
GOS72 & Indian Ocean               & --- \\
GOS73 & Indian Ocean               & --- \\
GOS74 & Indian Ocean               & --- \\
GOS75 & Indian Ocean               & --- \\
GOS76 & Coccos Keeling, Inside Lagoon & --- \\
GOS77 & Indian Ocean               & --- \\
GOS78 & Sargasso Sea               & --- \\
GOS79 & Sargasso Sea               & --- \\
GOS80 & Sargasso Sea               & --- \\
GOS81 & Sargasso Sea               & --- \\
\hline
\end{tabular}
\caption{Sample metadata: Part 2 of 2}
\label{tab:samples_part2}
\end{table}


\url{https://datacommons.cyverse.org/browse/iplant/home/shared/imicrobe/projects/26/samples}
